\chapter{The Software}

Two main types of software are required for the operation of this kind of aircraft. The flight controller runs on the on-board computer and is responsible for controlling the flight itself, and a ground station software, responsible for higher level commands and telemetry.
This chapters details these software used in the project.

\begin{figure}[h]
  \centering
  \begin{subfigure}{.5\textwidth}
    \centering
    \includegraphics[width=\linewidth]{figs/px4.png}
  \end{subfigure}%
  \begin{subfigure}{.5\textwidth}
    \centering
    \includegraphics[width=\linewidth]{figs/ardupilot.png}

  \end{subfigure}
  \caption{PX4 and Ardupilot.}
  \label{fig:motorcurves}
\end{figure}


\section{Flight Controller}
The flight controller software runs on the on-board computer, and is responsible for controlling the attitude, altitude, and position of the aircraft.
%
In order to achieve this, most flight controller boards come with internal sensors (accelerometer, gyroscope, magnetometer, barometer) and external ones (airspeed, magnetometer, GPS). 
%
The acquired data is used to estimate attitude and position, which is then corrected by the control loops.

The possible choices of flight controller software were ArduPilot and PX4.
%
Ardupilot is a community-drive software started by DYIDrones in 2009\cite{diydrones}.
%
It's GPL licensed, which means all changes made and commercialized must be open-sourced\cite{gplv3}.
%
ArduPilot is a mature software, with a large community of users and testers.
%
%https://github.com/PX4/Firmware/commits/master?after=3e0b8b7016556224aedb54f51d769d3288f0920e+24770

PX4 was also developed since 2008[\cite{waybackmachine}, mostly by the Computer Vision and Geometry Lab of ETH Zurich (Swiss Federal Institute of Technology)\cite{computervision} and the Autonomous Systems Lab\cite{autonomouslab} under a more permissive BSD 3-clause license\cite{bsd}.
%
For a while both projects worked closely, the Pixhawk Flight Controller board is a result of this interaction.
%
Both also joined Dronecode\cite{dronecode}, a Linux Foundation\cite{linuxfoundation} initiative started in 2004 as an attempt to grow the UAV ecosystem and reach larger companies.
%
%
Dronecode however, evolved into, according to the Ardupilot Dev Team, a flawed model.
%
It was required for all projects to hand over all domains, accounts, and trademarks to their control.
%
The project are also directed by the so-called "Platinum Members", which means the development would not be in control of the community anymore.
%
By September 9, 2016, a letter was released stating that Ardupilot was leaving Dronecode, and explaining why\cite{letter}.

The Ardupilot code was chosen for this project due to their larger openness and community.

